\SourceLeaf{161}
der Verpflichtungen der Kasse gegen alle künftigen Interessenten, oder der Verpflichtungen der dritten Rubrik dar:
\begin{center}
\begin{tabular}{l}
zu 42957 Thalern bei \(3\tfrac{1}{2}\) Proc. Zinsfuss\\
zu 32281 Thalern bei 4 Proc. Zinsfuss.
\end{tabular}
\end{center}
\noindent\hspace*{\parindent}Wollte man aber die jährlichen Beitritte durchschnittlich zu 3 veranschlagen, so würde man für diese Rubrik, je nach dem Zinsfuss, 8591 oder 6456 Thaler \textit{mehr} ansetzen müssen.

Fassen wir endlich die Resultate der bisherigen Anschläge zusammen, so ergibt sich aus den drei Rubriken
\begin{center}
\footnotesize\setlength{\tabcolsep}{3pt}
\begin{tabular}{l r|r}
&\(3\tfrac{1}{2}\) Proc.&4 Proc.\\
\shortstack[l]{Verpflichtungen gegen die jetzigen\\Witwen . . . . . . . . . . . . . . . . .}&46529 Thl.&44582 Thl.\\
\shortstack[l]{Verpflichtungen gegen die Hinterbliebenen\\der jetzigen Interessenten}&60724 ,, &54014 ,, \\
\shortstack[l]{Verpflichtungen gegen die künftigen\\Interessenten . . . . . . . . . . . . .}&42957 ,, &32281 ,, \\\cline{2-3}
\shortstack[l]{Totalsumme aller Verpflichtungen\\ . . . . . . . . . . . . . . . . . . . . . . . .}&150210 Thl.&130877 Thl.
\end{tabular}
\end{center}
\noindent\hspace*{\parindent}Von Nebenkosten will ich hier nur folgende Rubriken nach dem Durchschnitt der letzten 15 Jahre in Rechnung bringen (alles auf Gold reducirt):
\begin{center}
\small\setlength{\tabcolsep}{3pt}
\begin{tabular}{l r}
Bau- und Reparaturkosten, jährlich . . . . .&43 Rthl.\\
Processkosten nach Abzug der Erstattungen&39 ,, \\
Kosten der Rechnungsführung . . . . . . . . .&124 ,, \\\cline{2-2}
\multicolumn{1}{r}{Zusammen . . . . . .}&206 Rthl.
\end{tabular}
\end{center}
\noindent\hspace*{\parindent}Die Bau- und Reparaturkosten sind vielleicht nach diesem Durchschnitt zu niedrig angeschlagen. In dem oben S. [132] angeführten P. M. werden sie für einen Zeitraum von 24 Jahren (1768—1792) zu 1844 Thaler Cassenmünze (1975 Thl. 17 Ggr. Gold) angesetzt, was also einen jährlichen Durchschnittsbetrag von 82 Thalern ergeben würde: ich will jedoch bei obigem Mittel stehen bleiben. Diese 206 Thaler jährlicher Ausgabe repräsentiren ein Kapitalbedürfniss von 5886 oder von 5150 Thalern je nach dem Zinsfuss.

Es darf nicht übersehen werden, dass das bei allen Rechnungen zum Grunde liegende Interusurienprincip nur insoweit gültig ist, als alle Zinsen zu rechter Zeit eingehen und sofort wieder benutzt werden. In der Praxis ist diess aber nicht auszuführen, sondern es ist immer ein gewisser bald kleinerer bald grösserer Theil des Vermögens müssig. Man könnte den mittlern Betrag dieser Geldsumme zur Abkürzung Betriebsfonds nennen, obwohl diese Benennung genau zu reden nur auf denjenigen Theil des baaren Geldwerths passt, der bereit gehalten werden muss, damit die Kasse ihren eignen Zahlungsverpflichtungen pünktlich genügen könne. Es muss aber dazu auch eingerechnet werden das zeitweilig mussig liegende Geld, wenn sich nicht gleich Gelegenheit zu sicherer und vortheilhafter Belegung darbietet, und die ausstehenden Rückstände. Wie viel nun diess zusammen nach wirklichem mittlern Durchschnitte beträgt, lässt sich aus den Jahresrechnungen nur unvollkommen entnehmen, da dieselben nur angeben, \textit{was} im Laufe des Jahres eingenommen und verausgabt ist, nicht aber, an welchem Datum (dass diess \textit{zum Theil} aus den \textit{Belegen}, und auch aus diesen nicht ohne Ungewissheit, ergänzt werden \textit{könnte}, kommt hier nicht in Betracht). Ich habe indessen diesen Betrag aus den Rechnungsabschlüssen der letzten 11 Jahre\footnote{Früher stellten die Rechnungen den Vermögensstatus gar nicht auf.} so gut es geschehen kann abzuleiten gesucht, und danach gefunden: 2636 Thaler als mittlern Betrag des nicht zinstragenden Vermögenstheils.

Ich lasse nun die Vereinigung dieser verschiedenen Artikel hier folgen:
